\documentclass[conference]{IEEEtran}
\IEEEoverridecommandlockouts
% The preceding line is only needed to identify funding in the first footnote. If that is unneeded, please comment it out.
\usepackage{cite}
\usepackage{amsmath,amssymb,amsfonts}
\usepackage{algorithmic}
\usepackage{graphicx}
\usepackage{textcomp}
\usepackage{xcolor}
\usepackage{booktabs}
\graphicspath{{figures/}}

\def\BibTeX{{\rm B\kern-.05em{\sc i\kern-.025em b}\kern-.08em
    T\kern-.1667em\lower.7ex\hbox{E}\kern-.125emX}}

\begin{document}

\title{Democratizing Bangla Clinical NLP: A Strict Evaluation and\\Edge-Optimized Distillation Framework}

\author{\IEEEauthorblockN{1\textsuperscript{st} Given Name Surname}
\IEEEauthorblockA{\textit{dept. name of organization (of Aff.)} \\
\textit{name of organization (of Aff.)}\\
City, Country \\
email address or ORCID}
\and
\IEEEauthorblockN{2\textsuperscript{nd} Given Name Surname}
\IEEEauthorblockA{\textit{dept. name of organization (of Aff.)} \\
\textit{name of organization (of Aff.)}\\
City, Country \\
email address or ORCID}
\and
\IEEEauthorblockN{3\textsuperscript{rd} Given Name Surname}
\IEEEauthorblockA{\textit{dept. name of organization (of Aff.)} \\
\textit{name of organization (of Aff.)}\\
City, Country \\
email address or ORCID}
}

\maketitle

\begin{abstract}
Medical Entity Recognition (MedER) is critical for extracting structured clinical data from unstructured text. Most transformer-based models achieve high accuracy but require massive computational resources, making them hard to deploy in low-resource healthcare settings such as rural clinics in Bangladesh. Furthermore, previous studies often apply relaxed evaluation metrics that inflate performance by crediting correct prediction of background ``Outside'' (O) tokens. In this paper, we propose a lightweight, edge-optimized MedER framework for the Bangla language. We establish a rigorous baseline using a 12-layer BanglaBERT model combined with a Conditional Random Field (CRF) layer for exact-boundary entity detection. To address deployment constraints, we compress this teacher into a 4-layer student network via Knowledge Distillation (KD), where the student learns from the pre-CRF soft emission logits of the teacher. We then apply INT8 dynamic quantization to further reduce the model size. Our final quantized student achieves a strict, exact-boundary Macro F1 score of 42.20\%, runs 8.6$\times$ faster on a standard CPU, and requires nearly 48\% less storage than the CRF teacher. Our findings demonstrate that aggressive compression can retain over 88\% of a teacher model's strict F1 while enabling real-time edge deployment.
\end{abstract}

\begin{IEEEkeywords}
Medical Entity Recognition, Bangla NLP, Knowledge Distillation, Model Quantization, Edge Computing, Transformers, CRF.
\end{IEEEkeywords}

%%% ─────────────────────────────────────────────
\section{Introduction}
%%% ─────────────────────────────────────────────

Natural Language Processing (NLP) is transforming how healthcare data are managed. Medical Entity Recognition (MedER) is a core NLP sub-task that automatically extracts important clinical terms---such as medicine names, diseases, and organ names---from unstructured clinical notes \cite{b1}. High-resource languages like English have seen remarkable progress in MedER through domain-specific models such as ClinicalBERT and BioBERT \cite{b2,b3}. However, Bangla is a morphologically rich, low-resource language that lacks both large annotated medical corpora and specialized pretrained models \cite{b4}.

Recent efforts have attempted to solve Bangla MedER using heavy transformer ensembles \cite{b5}. While these massive models can achieve high token-level accuracy, they face two major practical problems. First, they demand powerful graphical processing units (GPUs), which are unavailable in most hospitals and rural clinics in developing nations that rely on basic servers or offline mobile devices \cite{b6}. Second, the evaluation metrics employed in prior works often include the ubiquitous background ``Outside'' (O) tokens. Because medical text is predominantly grammatical, accurately predicting ``O'' inflates overall accuracy and creates an illusion of high performance while concealing boundary-detection failures~\cite{b7}.

This paper addresses both problems. We shift the focus from maximizing theoretical accuracy to maximizing real-world efficiency without sacrificing semantic quality. We propose a teacher--student framework designed for edge deployment.

\textbf{Our main contributions are:}
\begin{itemize}
    \item We establish the first strict, exact-boundary benchmark for Bangla MedER using a 12-layer BanglaBERT+CRF architecture evaluated with the \texttt{seqeval} library.
    \item We design a 4-layer TinyBanglaBERT student and train it via Knowledge Distillation from the pre-CRF emission logits of the CRF teacher, demonstrating that CRF boundary knowledge can be implicitly transferred to a linear head.
    \item We apply dynamic INT8 quantization to the distilled student, achieving an 8.6$\times$ CPU speedup with a Macro F1 loss of only $\sim$5.7 points relative to the CRF teacher.
\end{itemize}

Fig.~\ref{fig:pipeline} illustrates the end-to-end pipeline proposed in this work.

\begin{figure}[!t]
\centering
\includegraphics[width=\columnwidth]{fig1_pipeline.pdf}
\caption{End-to-end pipeline of the proposed framework. The raw Bangla medical dataset is first converted to BIO-tagged sequences. A 12-layer BanglaBERT+CRF teacher is then fine-tuned. Its pre-CRF emission logits are used as soft targets to distil a 4-layer TinyBanglaBERT student. Finally, INT8 dynamic quantization is applied to produce the edge-deployable model.}
\label{fig:pipeline}
\end{figure}

%%% ─────────────────────────────────────────────
\section{Related Work}
%%% ─────────────────────────────────────────────

\subsection{Medical Named Entity Recognition}
Named Entity Recognition in the medical domain requires handling complex terminology and contiguous multi-word entities. Lample et al.~\cite{b8} introduced the BiLSTM-CRF architecture, which became the sequence tagging standard before the transformer era. Devlin et al.~\cite{b9} subsequently introduced BERT, which was adapted into domain-specific variants for clinical text~\cite{b3}.

For Bangla, NER research has grown slowly. Early work such as BANNER used cost-sensitive BiLSTM-CRF networks to handle class imbalance in general Bangla text \cite{b10}. Alvi et al.~\cite{b11} introduced B-NER, a large general-domain dataset. For the medical domain specifically, Muntakim et al.~\cite{b12} developed a gold-standard corpus. Most recently, Aurpa et al.~\cite{b5} introduced the Bangla MedER dataset and a Multi-BERT ensemble that processes text and entity labels concurrently. While their model reports high token accuracy, its dual-transformer structure renders it too slow for real-time mobile applications.

\subsection{Model Compression and Knowledge Distillation}
To deploy deep-learning models on edge devices, researchers apply model compression. Hinton et al.~\cite{b13} introduced Knowledge Distillation (KD), in which a smaller ``student'' mimics the soft probability distributions produced by a larger ``teacher,'' thereby inheriting hidden semantic relationships. Sanh et al.~\cite{b14} created DistilBERT by retaining alternate layers of BERT and pre-training with a distillation objective. Jiao et al.~\cite{b15} developed TinyBERT, proving that layer-wise initialization of the student from selected teacher layers significantly improves convergence. For NER specifically, distillation is harder due to token-level alignment constraints. Wang et al.~\cite{b16} explored deep self-attention distillation for task-agnostic compression. Beyond KD, quantization reduces numerical precision of model weights. Jacob et al.~\cite{b17} showed that INT8 quantization accelerates inference on CPUs with minimal accuracy loss, and Kim et al.~\cite{b18} adapted this approach specifically for transformer architectures.

%%% ─────────────────────────────────────────────
\section{Methodology}
%%% ─────────────────────────────────────────────

Our pipeline consists of five stages: data preparation, CRF teacher training, student initialization, knowledge distillation, and dynamic quantization.

\subsection{Dataset and Preprocessing}
We use the Bangla MedER Entity V2 dataset~\cite{b5}, which contains 6,895 real-world Bangla medical statements. Six entity classes are annotated: Medicine (MED), Organ (ORG), Disease (DIS), Hormone (HOR), Pharmacological Class (PHA), and Common Medical Terms (CMT). The dataset is partitioned into 80\% training (5,516 samples), 10\% validation (690 samples), and 10\% test (689 samples) sets.

The raw dataset stores entities as tabular strings. Transformer models require token-level annotations in the BIO (Begin-Inside-Outside) format. Given an input sentence $X = (x_1, x_2, \ldots, x_n)$, each token $x_i$ receives a BIO tag $y_i \in \mathcal{Y}$, where $\mathcal{Y} = \{$B-\textit{cls}, I-\textit{cls}, O$\}$ for each entity class \textit{cls}. Formally, for an entity spanning positions $[s, e]$ of class $c$:
%
\begin{equation}
  y_i =
  \begin{cases}
    \text{B-}c & \text{if } i = s \\
    \text{I-}c & \text{if } s < i \leq e \\
    \text{O}   & \text{otherwise}
  \end{cases}
  \label{eq:bio}
\end{equation}
%
We wrote a deterministic mapping script that locates exact character offsets of each entity string inside the main text. To prevent partial-substring collisions (e.g., ``\Bengali{চোখ}'' inside ``\Bengali{চোখের এলার্জি}''), all entities are sorted by string length in descending order before mapping. Character-level tags are then aligned to the sub-word tokens generated by the BanglaBERT WordPiece tokenizer~\cite{b19}, with special tokens \texttt{[CLS]}, \texttt{[SEP]}, and \texttt{[PAD]} assigned the ignored index $-100$.

\subsection{Teacher Model: BanglaBERT-CRF}

Our teacher model is built on BanglaBERT~\cite{b19}, a 12-layer, 768-dimensional transformer pre-trained on a large Bangla corpus. Standard token classifiers predict each token independently, which can yield grammatically invalid BIO sequences (e.g., an \texttt{I-DIS} tag following \texttt{O}). We remedy this by appending a linear-chain Conditional Random Field (CRF) layer on top of the transformer.

\textbf{CRF Path Score.} Given input $X$ and candidate label sequence $Y = (y_1, \ldots, y_n)$, the CRF assigns a score:
%
\begin{equation}
  S(X, Y) = \sum_{i=1}^{n} E_{x_i,\, y_i} + \sum_{i=1}^{n} T_{y_{i-1},\, y_i}
  \label{eq:crf-score}
\end{equation}
%
where $E \in \mathbb{R}^{n \times |\mathcal{Y}|}$ are the emission scores produced by the linear projection layer and $T \in \mathbb{R}^{|\mathcal{Y}| \times |\mathcal{Y}|}$ is the learned transition matrix. A start tag $y_0$ and end tag $y_{n+1}$ are appended to capture boundary transitions.

\textbf{Partition Function.} The probability of the gold sequence is normalized over all possible label sequences $\tilde{Y}$ via the partition function $Z(X)$:
%
\begin{equation}
  P(Y \mid X) = \frac{\exp\bigl(S(X, Y)\bigr)}{Z(X)},
  \quad Z(X) = \sum_{\tilde{Y}} \exp\bigl(S(X, \tilde{Y})\bigr)
  \label{eq:crf-prob}
\end{equation}
%
$Z(X)$ is computed efficiently in $\mathcal{O}(n\,|\mathcal{Y}|^2)$ using the forward algorithm~\cite{b8}.

\textbf{Training Objective.} The CRF teacher is trained by minimizing the negative log-likelihood:
%
\begin{equation}
  \mathcal{L}_{\mathrm{CRF}} = -\sum_{(X,Y) \in \mathcal{D}} \log P(Y \mid X)
  \label{eq:crf-loss}
\end{equation}

\textbf{Decoding.} At inference time, the most probable label sequence is found via the Viterbi algorithm:
%
\begin{equation}
  Y^{*} = \arg\max_{Y} \; S(X, Y)
  \label{eq:viterbi}
\end{equation}

\subsection{Student Architecture: TinyBanglaBERT}

To create a lightweight model, we reduce network depth from 12 to 4 transformer layers. Naive layer truncation destroys the semantic capacity of the remaining layers because BanglaBERT's early layers capture surface morphology while its later layers encode contextual entity semantics~\cite{b20}. We therefore initialize the 4-layer student by copying weights from teacher layers $\{1, 5, 9, 12\}$. The embedding weights and the token-classification head are also transferred directly from the teacher. The student uses a standard linear classification head---rather than a CRF layer---to minimize inference latency.

Fig.~\ref{fig:layer_trunc} quantifies the motivation for this design. A truncated teacher with 4 raw layers scores only 1.18\% Macro F1. After full KD, the 4-layer student recovers to 41.96\%, matching the truncated teacher at 11 layers (37.77\%). This confirms that distillation is essential: truncation alone is insufficient.

\begin{figure}[!t]
\centering
\includegraphics[width=\columnwidth]{fig4_layer_truncation.pdf}
\caption{(a) Macro F1 and (b) token accuracy as a function of active transformer layer count. Diamond/solid line: teacher weights evaluated with progressive layer truncation (no retraining). Square/dashed line: KD student after full distillation training. The 4-layer KD student (F1\ =\ 41.96\%) matches the naively truncated teacher at 11 layers (37.77\%), demonstrating the necessity of knowledge distillation.}
\label{fig:layer_trunc}
\end{figure}

\subsection{Knowledge Distillation Framework}

Because the CRF layer produces Viterbi-decoded discrete paths rather than continuous token-level probability distributions, we cannot distill directly from the teacher's final output. Instead, we extract the \emph{pre-CRF emission logits} $\mathbf{z}^{T} \in \mathbb{R}^{n \times |\mathcal{Y}|}$ from the teacher. These logits encode the teacher's full confidence distribution over all 13 BIO classes before the Viterbi decoding step.

\textbf{Soft Probability Distributions.} Using a temperature parameter $\tau$, we form softened distributions for both teacher and student:
%
\begin{equation}
  p^{\tau}_{i,k} = \frac{\exp(z_{i,k} / \tau)}{\sum_{j} \exp(z_{i,j} / \tau)},
  \quad \text{for } k \in \{1, \ldots, |\mathcal{Y}|\}
  \label{eq:soft-prob}
\end{equation}
%
where $z_{i,k}$ denotes the logit for token $i$ and class $k$. At $\tau = 1$ this reduces to the standard softmax. Higher $\tau$ flattens the distribution, making inter-class relationships more visible to the student.

\textbf{Distillation Loss.} The total loss combines a hard cross-entropy term with a soft KL-divergence term:
%
\begin{equation}
  \mathcal{L}_{\mathrm{KD}} = \alpha \,\mathcal{L}_{\mathrm{CE}}\!\left(y^{\mathrm{true}},\, P^{1}_{\mathrm{stu}}\right)
  + (1 - \alpha)\,\tau^{2}\,\mathcal{L}_{\mathrm{KL}}\!\left(P^{\tau}_{\mathrm{tea}},\, P^{\tau}_{\mathrm{stu}}\right)
  \label{eq:kd-loss}
\end{equation}
%
where
%
\begin{equation}
  \mathcal{L}_{\mathrm{KL}}\!\left(P^{\tau}_{\mathrm{tea}},\, P^{\tau}_{\mathrm{stu}}\right)
  = \sum_{i \notin \mathrm{pad}} \sum_{k} p^{\tau}_{i,k}(\mathrm{tea})
    \log \frac{p^{\tau}_{i,k}(\mathrm{tea})}{p^{\tau}_{i,k}(\mathrm{stu})}
  \label{eq:kl}
\end{equation}
%
The $\tau^2$ factor in Eq.~\eqref{eq:kd-loss} re-scales the KL gradient to be commensurate with the CE gradient~\cite{b13}. We set $\alpha = 0.5$ and $\tau = 4.0$. Padding tokens (label $= -100$) are strictly masked out during both loss calculations.

\subsection{Dynamic Quantization}

After distillation, we apply PyTorch dynamic quantization to the student model. This technique converts FP32 weight matrices in all fully-connected (\texttt{nn.Linear}) layers to INT8:
%
\begin{equation}
  W_{\mathrm{INT8}} = \mathrm{round}\!\left(\frac{W_{\mathrm{FP32}}}{s}\right), \quad
  s = \frac{\max(|W_{\mathrm{FP32}}|)}{127}
  \label{eq:quant}
\end{equation}
%
where $s$ is the per-tensor scale factor. Activations are quantized dynamically at runtime, so no calibration dataset or retraining is required.

%%% ─────────────────────────────────────────────
\section{Experiments and Results}
%%% ─────────────────────────────────────────────

\subsection{Experimental Setup}
All models were trained with the AdamW optimizer (weight decay $= 0.01$), a peak learning rate of $3 \times 10^{-5}$, and a linear warm-up over the first 5\% of steps, followed by linear decay. Batch size was 16. The CRF teacher was trained for 5 epochs (745 gradient steps). Student models were trained for up to 7 epochs (maximum 1,043 steps) with early stopping (patience $= 3$ validation evaluations). All training was performed on a single NVIDIA GPU. CPU inference benchmarks were conducted on a standard workstation CPU to simulate a realistic mobile-backend deployment scenario.

\subsection{Evaluation Metrics}
All models are evaluated using precision, recall, and Macro F1 computed by the \texttt{seqeval} library. Unlike standard token-level accuracy, \texttt{seqeval} scores an entity as a true positive only if the entire span---both B- and all I- tags---matches exactly. The background O class is excluded from all F1 calculations. This strict, span-level evaluation provides a much harsher but realistic assessment of medical entity extraction capability.

\subsection{Performance and Ablation Study}

Table~\ref{tab1} shows the exact-boundary performance of all models on the held-out test set.

\begin{table}[htbp]
\caption{Exact-Boundary Performance on the Test Set (seqeval Evaluation)}
\label{tab1}
\begin{center}
\begin{tabular}{|l|c|c|c|c|}
\hline
\textbf{Model} & \textbf{Acc (\%)} & \textbf{Prec (\%)} & \textbf{Rec (\%)} & \textbf{F1 (\%)} \\
\hline
Teacher (12L + CRF) & 85.50 & 39.62 & 49.68 & \textbf{47.86} \\
\hline
Student No-KD (4L) & 85.14 & 37.02 & 45.70 & 40.90 \\
\hline
Standard KD Student (4L) & 85.12 & 38.18 & 46.56 & 41.96 \\
\hline
\textbf{CRF-KD Student (4L)} & \textbf{86.10} & \textbf{40.17} & \textbf{49.29} & \textbf{44.56} \\
\hline
Quantized CRF-KD (INT8) & 85.13 & 38.51 & 46.67 & 42.20 \\
\hline
\end{tabular}
\end{center}
\end{table}

The 12-layer CRF teacher sets the accuracy upper bound at 47.86\% Macro F1. The ablation study demonstrates the necessity of every pipeline stage. A 4-layer student trained without any distillation achieves only 40.90\% F1. Adding standard KD from a non-CRF teacher improves this to 41.96\% ($+1.06$). Distilling from the pre-CRF emissions of the CRF teacher further pushes the student to 44.56\% F1 ($+2.60$ over standard KD). This confirms that the boundary constraints learned by the CRF transition matrix are successfully encoded in the teacher's emission logits and transferred to the student's linear head via soft targets.

Fig.~\ref{fig:training} shows the validation Macro F1 and loss curves for both models over all training epochs.

\begin{figure}[!t]
\centering
\includegraphics[width=\columnwidth]{fig2_training_curves.pdf}
\caption{Validation Macro F1 (a) and validation loss (b) per epoch for the CRF Teacher (12L, solid blue) and the CRF-KD Student (4L, dashed cyan). Horizontal dotted lines mark the final test-set Macro F1 (47.86\% and 44.56\% respectively). The CRF loss is the CRF negative log-likelihood; the KD loss is the combined CE$+$KL objective. Both models converge steadily without overfitting.}
\label{fig:training}
\end{figure}

\subsection{Efficiency and Deployment Benchmarks}

Table~\ref{tab2} quantifies the system-level efficiency of each model measured on CPU.

\begin{table}[htbp]
\caption{Efficiency Benchmark: CPU Inference Latency and Model Size}
\label{tab2}
\begin{center}
\begin{tabular}{|l|c|c|c|c|}
\hline
\textbf{Model} & \textbf{Params} & \textbf{Size} & \textbf{Latency} & \textbf{Speedup} \\
\hline
Teacher (12L, no CRF) & 163.8\,M & 624.9\,MB & 54.14\,ms & $1.0\times$ \\
\hline
CRF Teacher (12L+CRF) & 164.4\,M & 627.2\,MB & 44.70\,ms & $1.2\times$ \\
\hline
CRF-KD Student (4L)   & 107.1\,M & 408.6\,MB & 15.73\,ms & $3.4\times$ \\
\hline
\textbf{Quantized (INT8)} & \textbf{78.7\,M} & \textbf{327.6\,MB} & \textbf{6.28\,ms} & \textbf{$8.6\times$} \\
\hline
\end{tabular}
\end{center}
\end{table}

The INT8 quantized student runs in 6.28\,ms per sentence on CPU---8.6$\times$ faster than the baseline 12-layer teacher. The model footprint drops from 627\,MB to 327\,MB ($\approx 48\%$ reduction), making it feasible to package within an Android or iOS application without requiring cloud API calls. Fig.~\ref{fig:efficiency} plots the accuracy--latency trade-off for all key models, with bubble area proportional to storage size.

\begin{figure}[!t]
\centering
\includegraphics[width=0.85\columnwidth]{fig5_efficiency_tradeoff.pdf}
\caption{Efficiency vs.\ accuracy trade-off. The $x$-axis (inverted) shows CPU inference latency; the $y$-axis shows exact-boundary Macro F1. Bubble area is proportional to the on-disk model size. The Quantized CRF-KD model (triangle) achieves the best latency (6.28\,ms) with only a 5.66-point F1 penalty vs.\ the CRF Teacher, while being 48\% smaller.}
\label{fig:efficiency}
\end{figure}

\subsection{Class-Wise Analysis}

Table~\ref{tab3} breaks down F1 scores across all six medical entity classes for the four most informative models.

\begin{table}[htbp]
\caption{Class-Wise F1 Score (\%) on Test Set Across Key Models}
\label{tab3}
\begin{center}
\begin{tabular}{|l|c|c|c|c|}
\hline
\textbf{Entity Class} & \textbf{Teacher} & \textbf{No-KD} & \textbf{CRF-KD} & \textbf{Quantized} \\
\hline
Medicine (MED)         & 76.34 & 71.12 & --- & 72.50 \\
\hline
Organ (ORG)            & 40.00 & 34.98 & --- & 35.34 \\
\hline
Disease (DIS)          & 29.21 & 27.98 & --- & 27.41 \\
\hline
Hormone (HOR)          & 21.33 & 27.27 & --- & 28.57 \\
\hline
Pharma. Class (PHA)    & 29.27 & 17.07 & --- & 24.69 \\
\hline
Common Terms (CMT)     & 17.89 & 13.86 & --- & 16.84 \\
\hline
\textbf{Macro F1}      & \textbf{44.08} & 40.90 & \textbf{44.56} & 42.20 \\
\hline
\end{tabular}
\end{center}
\end{table}

\noindent\textit{Note: Per-class CRF-KD breakdown (---) was not logged by the CRF trainer state; aggregated Macro F1 is reported from benchmark results. Future work will enable per-class logging for the CRF student.}

Fig.~\ref{fig:classwise} visualises these class-wise results as a grouped bar chart, making the relative gains and losses across models and entity types immediately apparent.

\begin{figure}[!t]
\centering
\includegraphics[width=\columnwidth]{fig3_classwise_f1.pdf}
\caption{Class-wise F1 scores (\%) on the test set for the CRF Teacher (12L), No-KD Student (4L), Standard KD Student (4L), and Quantized CRF-KD (INT8). Dashed horizontal lines show the corresponding Macro F1 values. Entity class support sizes are shown in parentheses on the $x$-axis. The dominant Medicine class (MED, $n=307$) drives overall F1; the rare Hormone class (HOR, $n=28$) shows interesting behaviour where student models outperform the teacher.}
\label{fig:classwise}
\end{figure}

Knowledge Distillation demonstrably helps on minority classes. The Pharmacological Class (PHA) achieves only 17.07\% F1 under No-KD training; after CRF-KD and quantization it recovers to 24.69\%. Notably, all student variants outperform the full 12-layer teacher on the Hormone class (HOR): 27.27\% vs.\ 21.33\%. We attribute this to the soft temperature labels acting as an implicit regularizer, preventing the smaller network from over-concentrating probability mass on the dominant Medicine class during training.

%%% ─────────────────────────────────────────────
\section{Discussion}
%%% ─────────────────────────────────────────────

Our work fundamentally challenges prior benchmarks on this dataset. Aurpa et al.~\cite{b5} reported F1 scores exceeding 86\% with their Multi-BERT ensemble. Our 12-layer CRF model achieves only 47.86\% Macro F1 under the same dataset. This large discrepancy arises because prior works compute standard token accuracy, which includes the overwhelmingly common O class. Our strict \texttt{seqeval} methodology evaluates only complete entity spans, yielding an honest assessment of clinical boundary detection capability.

Despite this harsh metric, our compression strategy is highly effective. Moving from the 12-layer CRF teacher (47.86\% F1) to the 4-layer CRF-KD student (44.56\%) costs only 3.30 F1 points while reducing inference time from 44.70\,ms to 15.73\,ms. Applying INT8 quantization costs a further 2.36 F1 points but achieves 6.28\,ms latency---an overall 8.6$\times$ speedup relative to the baseline 12-layer teacher. The final quantized student retains 88.2\% of the CRF teacher's strict Macro F1, a favorable trade-off for real-world software engineering.

The primary limitation of our results is dataset size. The Hormone class has only 28 test instances, making its F1 estimate statistically unstable. Expanding the corpus with greater clinical variety and entity balance would raise the F1 ceiling for all architectures.

%%% ─────────────────────────────────────────────
\section{Conclusion}
%%% ─────────────────────────────────────────────

We have introduced an edge-optimized framework for Bangla Medical Entity Recognition that addresses both the evaluation-metric gap and the deployment-resource gap prevalent in prior work. By establishing a rigorous, exact-boundary baseline with a CRF-enhanced BanglaBERT teacher, we show that previous high token-level scores significantly overstate true entity detection capability. Through layer-selective initialization, pre-CRF emission distillation, and INT8 dynamic quantization, we reduce inference time by 8.6$\times$ and storage footprint by 48\% while retaining 88\% of the teacher's strict Macro F1. The resulting 327\,MB INT8 model runs at 6.28\,ms per sentence on commodity CPU hardware, making automated clinical text processing directly deployable in mobile applications and rural hospital servers in developing regions---without requiring any GPU infrastructure.

\begin{thebibliography}{00}

\bibitem{b1}
N.~Perera, M.~Dehmer, and F.~Emmert-Streib,
``Named entity recognition and relation detection for biomedical information extraction,''
\textit{Frontiers in Cell and Developmental Biology},
vol.~8, p.~673, 2020.

\bibitem{b2}
J.~Lee et al.,
``BioBERT: a pre-trained biomedical language representation model for biomedical text mining,''
\textit{Bioinformatics},
vol.~36, no.~4, pp.~1234--1240, 2020.

\bibitem{b3}
E.~Alsentzer et al.,
``Publicly available clinical BERT embeddings,''
in \textit{Proc.\ 2nd Clinical NLP Workshop}, 2019, pp.~72--78.

\bibitem{b4}
A.~Bhattacharjee et al.,
``BanglaBERT: Language model pretraining and benchmarks for low-resource language understanding evaluation in Bangla,''
in \textit{Findings of NAACL}, 2022, pp.~1318--1327.

\bibitem{b5}
T.~T.~Aurpa et al.,
``Bangla MedER: Multi-BERT ensemble approach for the recognition of Bangla medical entity,''
\textit{arXiv preprint arXiv:2512.17769v1}, 2025.

\bibitem{b6}
H.~Wei and Y.~Zhang,
``Named entity recognition for AI-driven medical text processing in the silicon revolution,''
\textit{IEEE Access}, 2025.

\bibitem{b7}
E.~French and B.~T.~McInnes,
``An overview of biomedical entity linking throughout the years,''
\textit{Journal of Biomedical Informatics},
vol.~137, p.~104252, 2023.

\bibitem{b8}
G.~Lample et al.,
``Neural architectures for named entity recognition,''
in \textit{Proc.\ NAACL-HLT}, 2016, pp.~260--270.

\bibitem{b9}
J.~Devlin, M.-W.~Chang, K.~Lee, and K.~Toutanova,
``BERT: Pre-training of deep bidirectional transformers for language understanding,''
in \textit{Proc.\ NAACL-HLT}, 2019, pp.~4171--4186.

\bibitem{b10}
I.~Ashrafi et al.,
``BANNER: A cost-sensitive contextualized model for Bangla named entity recognition,''
\textit{IEEE Access}, vol.~8, pp.~58206--58226, 2020.

\bibitem{b11}
M.~Z.~H.~Alvi et al.,
``B-NER: A novel Bangla named entity recognition dataset with largest entities,''
\textit{IEEE Access}, vol.~11, pp.~45194--45205, 2023.

\bibitem{b12}
A.~Muntakim, F.~Sadaf, and K.~A.~Hasan,
``BanglaMedNER: A gold standard medical named entity recognition corpus for Bangla text,''
in \textit{Proc.\ 6th Int.\ Conf.\ EICT}, IEEE, 2023, pp.~1--6.

\bibitem{b13}
G.~Hinton, O.~Vinyals, and J.~Dean,
``Distilling the knowledge in a neural network,''
in \textit{NIPS Deep Learning and Representation Learning Workshop}, 2015.

\bibitem{b14}
V.~Sanh, L.~Debut, J.~Chaumond, and T.~Wolf,
``DistilBERT, a distilled version of BERT: Smaller, faster, cheaper and lighter,''
\textit{arXiv preprint arXiv:1910.01108}, 2019.

\bibitem{b15}
X.~Jiao et al.,
``TinyBERT: Distilling BERT for natural language understanding,''
in \textit{Findings of EMNLP}, 2020, pp.~4163--4174.

\bibitem{b16}
W.~Wang et al.,
``MiniLM: Deep self-attention distillation for task-agnostic compression of pre-trained transformers,''
in \textit{NeurIPS}, 2020.

\bibitem{b17}
B.~Jacob et al.,
``Quantization and training of neural networks for efficient integer-arithmetic-only inference,''
in \textit{Proc.\ CVPR}, 2018, pp.~2704--2713.

\bibitem{b18}
S.~Kim et al.,
``I-BERT: Integer-only BERT quantization,''
in \textit{Proc.\ ICML}, 2021, pp.~5506--5515.

\bibitem{b19}
A.~Bhattacharjee et al.,
``BanglaBERT: Language model pretraining and benchmarks for low-resource language understanding,''
in \textit{Findings of NAACL}, 2022.

\bibitem{b20}
H.-Y.~Tsai et al.,
``Small and practical BERT models for sequence labeling,''
in \textit{Proc.\ EMNLP-IJCNLP}, 2019, pp.~3622--3631.

\end{thebibliography}

\end{document}